% ============================================================
% VERSIÓN INDIVIDUAL — DANIEL DUARTE CORDERO
% Se rellenan únicamente afirmaciones sostenidas por la auditoría
% del proyecto y por las referencias leídas/asignadas a Daniel.
% Los TODO de integración corresponden a aportes del equipo.
% ============================================================
\section{Introducción}
\label{sec:introduccion}

\subsection{Contexto y relevancia}
\label{subsec:contexto}
Los juegos competitivos generan secuencias de estados en las que la
ventaja cambia conforme avanza la partida. Estimar la probabilidad de
victoria desde información disponible durante el juego es, por tanto,
un problema distinto de predecir el ganador únicamente antes del
inicio. En deportes electrónicos de tipo MOBA, la literatura ya trata
la predicción de resultado como una aplicación consolidada de
aprendizaje automático y emplea tanto variables previas a la partida
como información observada durante ella \cite{kamal_moba_2025}. Un
ejemplo directo es el trabajo de Tian et al., que realiza predicción en
tiempo real en \emph{Honor of Kings} a partir de secuencias de estados
y muestra que la importancia de la alineación, el oro, las bajas y las
torres cambia con el progreso del encuentro \cite{tian_hero_2022}.

La misma idea de asignar valor a estados intermedios aparece en juegos
de cartas. Miernik y Kowalski evolucionan funciones de evaluación para
\emph{Legends of Code and Magic}, aunque su salida es un puntaje
heurístico optimizado indirectamente por el rendimiento del agente, no
una probabilidad calibrada de victoria \cite{miernik_evolving_2022}. Más
cerca de nuestra formulación, el reto AAIA'17 de Hearthstone pidió
estimar la probabilidad de que un jugador ganara a partir de un estado
intra-partida y evaluó las soluciones mediante AUC
\cite{janusz_hearthstone_2017}. En consecuencia, la pregunta general de
predecir el desenlace desde estados intermedios no se presenta aquí
como inédita; el interés está en trasladarla rigurosamente al dominio
de Pokémon TCG y estudiar sus problemas específicos de representación,
fuga de información y generalización.

La generalización merece atención especial porque los entornos
competitivos cambian. Chitayat et al. muestran en Dota~2 que
representaciones ligadas a parámetros de diseño pueden conservar su
rendimiento al atravesar parches posteriores, mientras que
representaciones basadas solo en identidad son más frágiles
\cite{chitayat_meta_2023}. En Pokémon VGC, VGC-Bench separa explícitamente
el rendimiento sobre equipos vistos de la generalización a equipos no
vistos y muestra que entrenar con mayor diversidad de composiciones
mejora esa segunda capacidad \cite{angliss_vgcbench_2026}. Estos
antecedentes motivan que una evaluación de Pokémon TCG no se limite a
una partición aleatoria de partidas, sino que también examine cambios
temporales y composiciones no observadas durante el entrenamiento.

\subsection{Planteamiento del problema}
\label{subsec:problema}
El conjunto \emph{PTCG AI Battle Challenge Simulation Episodes} contiene
replays completos entre agentes artificiales. Una partida produce
muchos snapshots correlacionados, por lo que la unidad de observación
no coincide con la unidad estadísticamente independiente. La auditoría
preliminar del proyecto identificó 52,085 episodios independientes y
más de siete millones de snapshots útiles. Si estos snapshots se
particionaran al azar, estados de una misma partida podrían aparecer en
entrenamiento y prueba, produciendo una estimación artificialmente
optimista del desempeño. El problema metodológico central es entonces
estimar $P(\mathrm{win}\mid\mathrm{state})$ usando únicamente la
información observable hasta el instante de predicción y manteniendo
intactas las unidades independientes durante la evaluación.

Además, un buen resultado bajo partidas de distribución similar no
basta para demostrar que el modelo aprendió propiedades transferibles
del estado. La población auditada cambia en el tiempo en mazos, agentes
y versión del motor, y los mazos más frecuentes concentran una parte
importante de las observaciones. Por ello, el estudio debe distinguir
entre generalización a nuevas partidas, a periodos futuros y a
composiciones no vistas. Esta separación sigue la preocupación por
robustez entre parches de Chitayat et al. y por composiciones no vistas
de VGC-Bench \cite{chitayat_meta_2023,angliss_vgcbench_2026}.

\subsection{Pregunta de investigación}
\label{subsec:pregunta}
\begin{quote}
\textbf{PI:} ¿En qué medida las características observables del estado
de una partida de Pokémon TCG permiten predecir su resultado final
mediante modelos clásicos de aprendizaje automático, cómo cambia esa
capacidad predictiva conforme avanza la partida, y cuánto se conserva
al evaluar sobre partidas, mazos y periodos no vistos en el
entrenamiento?
\end{quote}

\noindent\textbf{Por qué es refutable.} La respuesta sería negativa o
quedaría sustancialmente debilitada si el intervalo de confianza del
ROC-AUC en la ventana temprana incluyera 0.5, si el desempeño no
mejorara al avanzar la partida, o si la señal desapareciera al evaluar
en el protocolo temporal o en el holdout de mazos/matchups.

\noindent\textbf{Por qué es medible.} La discriminación se medirá con
ROC-AUC y la calidad de la probabilidad con log loss, Brier score y
curvas de calibración. Las métricas se reportarán globalmente y por
ventana de turno. Las particiones y los intervalos de confianza se
construirán sobre episodios, aunque cada fila del dataset tabular sea un
snapshot.

\subsection{Hipótesis de trabajo}
\label{subsec:hipotesis}
\begin{quote}
\textbf{H1:} Las características estructuradas del estado contienen
señal predictiva sobre el resultado desde etapas tempranas
(ROC-AUC significativamente mayor que 0.5) y esa señal aumenta conforme
progresa la partida.
\end{quote}
\begin{quote}
\textbf{H2:} La señal temprana no se explica únicamente por el orden de
turno ni por la identidad del mazo.
\end{quote}

\noindent\textbf{Criterios de falsación.} H1 temprana se debilita si el
intervalo de confianza del AUC temprano incluye 0.5; H1 dinámica se
debilita si el AUC por ventana no aumenta una vez controlada la
longitud de partida. H2 se debilita si un baseline compuesto por
\texttt{is\_first\_player} e identidad de mazo iguala al modelo de estado
completo. Una caída hasta azar en el protocolo temporal se interpretará
como falta de generalización, no como un fallo de ejecución.

\subsection{Objetivos}
\label{subsec:objetivos}
\textbf{Objetivo general.} Desarrollar y evaluar modelos clásicos de
aprendizaje automático que estimen la probabilidad de victoria en
Pokémon TCG a partir de características tabulares derivadas
exclusivamente del estado observable, y caracterizar cómo evoluciona
esa capacidad predictiva a lo largo de la partida.

\noindent\textbf{Objetivos específicos.}
\begin{enumerate}
  \item Caracterizar el dataset y su calidad mediante EDA, diccionario
  de datos, análisis de faltantes, valores atípicos, desbalance y
  sesgos sobre la representación tabular final.
  \item Construir un pipeline reproducible desde replay JSON hasta una
  tabla de snapshots frescos con perspectiva normalizada.
  \item Diseñar y verificar controles de fuga mediante agrupación por
  \texttt{episode\_id}, lista negra de campos y pruebas automáticas.
  \item Comparar una escalera de baselines clásicos bajo protocolos
  agrupado, temporal y de robustez por mazo/matchup.
  \item Analizar discriminación, calibración, explicabilidad y
  generalización globalmente y por ventana de turno.
\end{enumerate}

\subsection{Contribución esperada}
\label{subsec:contribucion}
Se espera aportar una formulación reproducible de estimación de
probabilidad de victoria desde estados intermedios de Pokémon TCG, un
pipeline tabular con controles explícitos de fuga, y una comparación
entre modelos clásicos bajo distintos ejes de generalización. El
trabajo también hará explícita la diferencia entre discriminación y
calidad probabilística y evaluará si la información útil del estado
cambia conforme avanza la partida, sin prometer de antemano que exista
una mejora positiva.

\subsection{Estructura del documento}
\label{subsec:estructura}
La Sección~\ref{sec:relacionados} revisa la literatura en tres ejes y
cierra identificando las brechas. La Sección~\ref{sec:baseline} declara
el baseline de literatura y su comparabilidad. La
Sección~\ref{sec:metodologia} presenta la metodología propuesta. La
Sección~\ref{sec:analisistecnico} desarrolla el análisis técnico
específico del track. La Sección~\ref{sec:etica} documenta las
consideraciones éticas preliminares.
